\section{Carnap: comparison, relevance, and reliability proofs}
\label{sec:annex-p15-carnap-confirmation-relevance-and-reliability}

The definitions and standing in
Chapter~\ref{sec:p15-carnap-confirmation-relevance-and-reliability}
apply throughout. A finite system may contain arbitrary quantified
sentences and primitive relations; their ranges are finite sets of
complete state descriptions. No choice of \(c^*\) is made except in
the explicitly named examples. The proofs reconstruct the corresponding
1950 source claims independently; they do not assert that a later
family or probability-space existence theorem was present in that source.

\begin{annexitem}{annex:p15.carnap-comparative}{thm:p15-carnap-comparative}{The full finite comparison criterion}
\begin{proof}
Write \(A=R(eh)\), \(B=R(e\neg h)\),
\(C=R(e'h')\), and \(D=R(e'\neg h')\). The two conditional
denominators are positive. Cross multiplication and cancellation of
\(m(A)m(C)\) show that their comparison is equivalent to
\[
 m(A)m(D)\geq m(C)m(B).
\]
If \(B\) is empty or \(C\) is empty, the comparison holds for
every measure. Otherwise it holds for every positive state measure
exactly when \(C\subseteq A\) and \(B\subseteq D\).
Sufficiency follows by monotonicity of the two factors.

For necessity, if \(z\in C\setminus A\), choose
\(y\in B\). A probability supported on these one or two states
has zero mass on \(A\) and positive masses on \(C,B\); its
cross-multiplied comparison is strictly false. Similarly, if
\(z\in B\setminus D\), choose \(y\in C\). The supported
measure has zero mass on \(D\) and positive masses on \(B,C\),
and again the comparison is strictly false. In either construction,
mix the supported measure with a positive uniform measure on the
whole finite state space with weight \(\varepsilon>0\). All state
masses become positive. The cross-product difference is a polynomial
in \(\varepsilon\), still negative for sufficiently small positive
\(\varepsilon\) by continuity. Thus a regular countermeasure
exists, even though the limiting supported measure itself is not regular.

The logical condition \(C\subseteq A\) is
\(e'h'\models eh\). Under this condition,
\(B\subseteq D\) is equivalent to \(e\models h\lor e'\):
the latter puts \(B\) in \(R(e')\), and it cannot intersect
\(C\), since \(B\cap A=\varnothing\). Conversely
\(B\subseteq D\subseteq R(e')\) gives that implication.
The two empty-range cases are exactly \(e\models h\) and
\(e'\models\neg h'\). This proves D81-1's characterization
over all regular measures and hence its uniqueness as the most
comprehensive unanimous comparison.

Applying the characterization in both directions proves the Eq
assertion. Applying it in one direction and negating its converse proves
the stated Gr quantifiers. If neither direction holds, the just-proved
necessity constructs a regular measure violating each ordering, giving
the two opposite strict comparisons. For a common evidence,
\(C\subseteq A\) means \(eh'\models h\); the empty cases
are already included in this inclusion. Maximum and minimum follow
because positive mass on every admissible state makes a conditional
value one or zero precisely when its failure or success range is empty.

For the limit statement, every eventual finite comparison is an
inequality of real numbers. If both sequences converge, taking their
limits preserves the weak inequality. For nongeneral sentences, all
the relevant truth conditions use a fixed finite set of ground atoms.
Any chosen positive measure on the states involving those names can
be extended to larger systems by assigning all newly introduced ground
atoms independent fair truth values, including relation instances
involving new names. The extended measures are positive and fit under
restriction. Every sentence in the given comparison keeps its range
mass and confirmation value, so its infinite-language limit exists and
is that same value. In particular every finite countermeasure above
extends to a counterexample limit. Conversely any fitting regular
extension must obey an eventual finite unanimous ordering. This proves
the nongeneral characterization, not an arbitrary quantified-sentence
limit-existence or logical-transfer theorem. The source distinction is
T81-1 versus T81-2, pp.~440--441 \cite{Carnap1950Probability}.
\end{proof}
\end{annexitem}

\begin{annexitem}{annex:p15.carnap-comparative-calculus}{thm:p15-carnap-comparative-calculus}{Composition, complement, and predictable observations}
\begin{proof}
For finite systems, the preceding adequacy theorem reduces MC to a
comparison true for every regular measure. Real weak order is reflexive
and transitive, so MC is too. Mutual comparison gives an equivalence
relation. If a weak chain contains a Gr link, a measure making that
link strict makes the composite comparison strict; its reverse cannot
hold universally. This proves transitivity of Gr and its irreflexivity.
Complementation replaces each conditional value \(p\) by \(1-p\),
reversing the ordering.

For each regular measure put
\(a=c(h,e)\), \(b=c(j,eh)\),
\(a'=c(h',e')\), and \(b'=c(j',e'h')\).
The multiplication identity gives \(c(hj,e)=ab\) and
\(c(h'j',e')=a'b'\). If \(a\geq a'\) and \(b\geq b'\),
their nonnegative products have that ordering. The crossed premises
\(a\geq b'\), \(b\geq a'\) give the same result. For
cancellation, \(ab\geq a'b'\) and \(a'\geq a>0\) imply
\(b\geq(a'/a)b'\geq b'\). Exchanging the names of the two
left factors or the two right factors proves each of the four versions.
Consistency of the relevant conjunctive evidences guarantees the
positive factors needed for division. Since these conclusions hold for
every regular measure, the adequacy theorem turns them back into MC.

The disjunctive-evidence assertion must not be proved by taking an
arbitrary selected measure's two bounds: overlapping evidence need
not yield a convex combination of those two conditional values.
Instead use the logical criterion. If \(e\models h\), the
conclusion is immediate. Otherwise set \(A=R(eh)\),
\(B=R(e\neg h)\), and let
\(C_i=R(e'ih')\), \(D_i=R(e'i\neg h')\), with analogous
\(C_j,D_j\). Each premise says either that its \(C\)-range is
empty, or that \(C\subseteq A\) and \(B\subseteq D\).
If both success ranges are empty, their union is empty. Otherwise
their union is contained in \(A\), and the \(D\)-range of
at least one nonempty-success premise contains \(B\). Thus
\(C_i\cup C_j\subseteq A\) and
\(B\subseteq D_i\cup D_j\), giving MC for the union evidence.
This reproduces the content of T83-37 without reversing quantifiers.

For predictable \(i\), \(eh\models i\) gives
\(R(ehi)=R(eh)\). Hence
\(c(h,ei)=m(eh)/m(ei)=c(h,e)/c(i,e)\).
Both numerator and denominator are positive, and
\(0<c(i,e)\leq1\). Equality holds exactly when
\(c(i,e)=1\), which in a regular finite system is
\(e\models i\). For two predictable observations, dividing
the same positive numerator by their expectedness reverses ordering.
For two hypotheses predicting the same observation, division by its
common positive expectedness preserves ordering. These are the
qualified results of T85-10--12 \cite{Carnap1950Probability}.
\end{proof}
\end{annexitem}

\begin{annexitem}{annex:p15.carnap-relevance-core}{thm:p15-carnap-relevance-core}{The determinant, symmetry, and two additivities}
\begin{proof}
The four cells are disjoint and cover \(R(e)\). Their sums give
\(m(e)=a+b+c+d\), \(m(ei)=a+b\), and
\(m(eh)=a+c\). Substitute into the definition of \(r\):
\[
 a(a+b+c+d)-(a+c)(a+b)=ad-bc.
\]
When \(m(e),m(ei)>0\), subtract the two conditional ratios
over their common positive denominator to obtain the difference formula.
If \(ei\) is contradictory, then \(a=b=0\), \(r=0\),
and irrelevance holds by D65-1d's convention without forming a ratio.
If it is consistent in a finite regular system, both denominators are
positive and the sign correspondence follows. This also covers a
contradictory prior \(e\) by its convention, although no conditional
value is assigned there.

Exchanging observation and hypothesis interchanges \(b,c\) and
leaves \(ad-bc\) unchanged. Negating the observation interchanges
\((a,b)\) with \((c,d)\), reversing the sign; negating the
hypothesis interchanges \((a,c)\) with \((b,d)\), doing the
same. With \(m(e)>0\), division of the defining identity by
\(m(e)^2\) gives
\[
 r/m(e)^2=c(ih,e)-c(i,e)c(h,e),
\]
proving the special multiplication assertion, including zero factors.

Apply inclusion--exclusion to each of
\(m(eh(i\lor j))\) and \(m(e(i\lor j))\). Substitution
into \(r\)'s definition gives its inclusion--exclusion identity.
If \(m(eij)=0\), the smaller intersection with \(h\) also
has mass zero, so its signed relevance vanishes. If
\(m(e\neg i\neg j)=0\), complementation gives
\(r(i\lor j,h,e)=0\); rearrange inclusion--exclusion to obtain
conjunction additivity. Induction proves the multiple disjoint-union
identity, and complementation proves the analogous multiple conjunction
identity when the negated components are pairwise mass-exclusive.
Symmetry transfers all these results to the hypothesis argument. Only
finite additivity and the existing indicated measure values were used
in this algebra, so it remains valid beyond finite systems on the stated
domain. This reconstructs T65-4--6, T67-1--10, T68-1--4 and their
T70 counterparts \cite{Carnap1950Probability}.
\end{proof}
\end{annexitem}

\begin{annexitem}{annex:p15.carnap-relevance-quotients}{thm:p15-carnap-relevance-quotients}{Symmetric quotients and cumulative dependence}
\begin{proof}
The multiplication identity gives
\(Q_e(h,i)=c(hi,e)/(c(h,e)c(i,e))\), proving symmetry on its
domain. Replacing \(c(hi,e)\) with the determinant identity
proves the formula involving \(r\). The multiple expression is
permutation-invariant and its product formula is immediate from its
definition. Grouping its first event against the conjunction of the
remaining events gives the recursive form used in D66-1.

Put \(I=\bigwedge_r i_r\), and write
\(L_r=c(i_r,eh)/c(i_r,eh')\). Bayes' ratio form gives
\(O_r=O_0L_r\) and
\(c(h,eI)/c(h',eI)=O_0c(I,eh)/c(I,eh')\), on the stipulated
positive denominators. The two multiple quotient identities yield
\[
 \frac{c(I,eh)}{c(I,eh')}=
 \frac{Q_{eh}(i_1,\ldots,i_n)}{Q_{eh'}(i_1,\ldots,i_n)}
 \prod_r L_r.
\]
Substitute \(L_r=O_r/O_0\) and multiply by \(O_0\).
This gives the asserted cumulative formula, including \(n=1\)
when each one-argument dependence quotient is one. None of the
dependence quotients has been assumed to equal one. This is the
qualified Johnson--Keynes cumulative method of §66
\cite{Carnap1950Probability}.
\end{proof}
\end{annexitem}
\begin{annexitem}{annex:p15.carnap-relevance-sign-cases}{thm:p15-carnap-relevance-sign-cases}{All sign patterns and the original countercases}
\begin{proof}
The four observation cells partition the state space. Their signed
relevances sum to \(r(\top,h,e)=0\), by disjoint additivity.
The cells forming \(i,j,i\lor j\) give the three sums in the
statement. Conversely, let \(x_1+\cdots+x_4=0\). Use eight
states distinguished by the truth values of \(i,j,h\), taking
\(e=\top\). In observation cell \(l\), assign its \(h\)-true
state mass \(1/8+\varepsilon x_l\) and its \(h\)-false state
mass \(1/8-\varepsilon x_l\). Choose \(\varepsilon>0\) small
enough that all masses are positive. Each cell has mass \(1/4\),
the hypothesis has mass \(1/2\), and its relevance to cell \(l\)
is exactly \(\varepsilon x_l\). Thus every feasible sign tuple
is realized in a regular language with three independent ground atoms.

To count the tuples, nonzero vectors must have both positive and
negative coordinates. If all four coordinates are nonzero, eight
patterns have exactly one positive or exactly one negative coordinate;
their two additional sum signs are forced. The six two-positive,
two-negative patterns are
\[
 ++--,\quad +-+-,\quad +--+,\quad -++-,\quad -+-+,\quad --++.
\]
For the first, second, fifth, and sixth, one of the two sum signs is
forced and the other has all three possibilities. For the third and
fourth, the two signs vary independently, giving nine possibilities
each. Independence here can be seen by fixing positive \(x_1\)
and \(x_4\) and allowing their total to be split between the two
negative coordinates: \(x_1+x_2\) and \(x_1+x_3\) can independently
be positive, negative or zero by choosing sufficiently large or small
\(x_1\) relative to the two negative magnitudes; their sum is
\(x_1-x_4\), which is unrestricted. Explicitly, for prescribed
sum values \(u,v\), choose
\(x_1>\max(0,u,v,u+v)\), set
\(x_2=u-x_1\), \(x_3=v-x_1\), and
\(x_4=x_1-u-v\). This realizes the \(+--+\) pattern for any
signs of \(u,v\); negate it for \(-++-\).
The count is therefore \(8+4\cdot3+2\cdot9=38\).
With exactly one zero, choose its position in four ways and a mixed
sign assignment to the other three in six ways. Their additional sum
signs are forced by zero-sum, giving twenty-four. With two zeros,
choose the two nonzero positions in six ways; their signs are opposite
and can be assigned in two ways, giving twelve. Exactly three zeros
forces four zeros. Including the all-zero tuple gives
\(38+24+12+1=75\). This is the complete T69-1 classification.

Put \(u=x_1+x_2\), \(v=x_1+x_3\),
\(a=x_1\), and \(b=-x_4\). Then \(a+b=u+v\).
If \(u,v\geq0\) and at least one is positive, at least one of
\(a,b\) is positive. Every other allowed pair of signs can be
realized by choosing real \(a,b\) of those signs with positive
sum and splitting their sum into the stipulated \(u,v\).
For \(u>0,v<0\), any chosen \(a,b\) is possible: take
\(u=1+\max(a+b,0)\), \(v=a+b-u\), and recover
\(x_2=u-a,x_3=v-a,x_4=-b\). If \(u=v=0\), then
\(b=-a\). Negation gives the negative cases. A logical exclusion
or implication simply removes its impossible observation cell; all
identities remain valid. Under an implication, all nine signs for
\(u,v\) remain realizable by taking \((u,0,v-u,-v)\) and
using three nonempty observation cells, each with both \(h\)-values
and a sufficiently small perturbation of equal conditional masses.
Outside the consistent evidence imposing that implication, give the
remaining states positive mass. Symmetry of \(r\) proves the
two-hypothesis version. These calculations establish T69-2--5 and
T71-2--5's sign content without a selected inductive family.

For the tournament, conditional outcome masses are \(1/10\).
The local, junior, and male sets have cardinality five; local males
and junior males have cardinality three; local junior males cardinality
one; local juniors cardinality three. Division gives exactly
\(1/2,3/5,1/3,3/10,1/5\) as stated. Visiting and senior sets
each contain three women; their union contains seven players and
three women, giving \(3/7\). Its prior probability is \(7/10\),
and conditional on a woman its probability is \(3/5\). Complementing
the gender event gives the opposite changes. These ten outcomes can
be embedded in a regular state space of sixteen states: take the
prior evidence to select ten, use uniform positive masses on all
sixteen, and define the tournament events as the appropriate unions.
Thus the examples respect regularity rather than introducing zero
state masses tacitly. Source locators: §§69 and 71, pp.~382--397
\cite{Carnap1950Probability}.
\end{proof}
\end{annexitem}

\begin{annexitem}{annex:p15.carnap-relevance-analysis}{thm:p15-carnap-relevance-analysis}{Both atomic analyses and all stronger relevance conditions}
\begin{proof}
The range of \(i\) is the disjoint union of its single-state ranges,
so additivity gives its first sum. A state outside \(R(e)\) has
zero relevant intersection. A state satisfying \(eh\) has
\(r(z,h,e)=m(z)(m(e)-m(eh))=m(z)m(e\neg h)\);
one satisfying \(e\neg h\) has
\(r(z,h,e)=-m(z)m(eh)\). These prove the atomic formulas.
The sum over every state is zero; negation reverses each term's sign.
Therefore the sum of \(r(\neg z,h,e)\) over states falsifying
\(i\) equals the sum over states satisfying it. Equivalently,
\(i\) excludes precisely those states, and is the conjunction of
their negations. These negations are pairwise exhaustive in disjunction,
so conjunction additivity gives the same result. Complementation of
ranges also proves that strengthening a sentence enlarges its excluded
content. The empty range and empty content cases use empty sums zero
and the empty conjunction true.

If \(i\) is positive, its determinant forces \(a,d>0\).
Any state in \(ei\neg h\) is a stronger negative observation
by the atomic formula. Thus extreme positivity requires \(b=0\).
Conversely with \(b=0,a,d>0\), every state of every stronger
compatible observation lies in \(eh\), so no such observation
is negative. For complete positivity, a state in \(e\neg i h\)
has a negative content-negation implied by \(ei\), so \(c=0\)
is necessary. If \(c=0,a,d>0\), every excluded state of any
weaker observation is outside \(e\) or in \(e\neg h\);
its negation has nonnegative relevance. Their sum proves sufficiency.
Changing \(h\) to \(\neg h\) gives the negative criteria.

For extreme irrelevance, every state in \(ei\) must have zero
relevance. If \(ei\) is empty this is automatic. Otherwise the
atomic formulas and positivity of its state masses force either
\(m(eh)=0\) or \(m(e\neg h)=0\). These conditions also
suffice, since the prior evidence already decides the hypothesis.
The content analysis gives the same argument with \(e\neg i\)
for complete irrelevance. Positivity of state masses turns each zero
mass into logical contradiction, proving measure-independence.
The criteria immediately give the strengthened-observation and
weakened-observation transmission statements, with their compatibility
or nonentailment restrictions.

For a positive determinant, \(a,d>0\); hence extreme and complete
are respectively \(b=0,c=0\). If both are nonzero, regularity
allows their relative weights to vary independently of \(a,d\),
so the determinant can have any sign. If at least one is structurally
zero, it is positive for every regular measure. Reversing the products
gives the negative cases. A determinant is zero for every regular
measure exactly when both products are structurally zero; otherwise
positive weights can make it nonzero. This supplies all sixteen
zero-support patterns of T72-8, with its three determinant subdivisions
when all four cells are present, and the subdivisions in T76-1--4.
Negating \(i\) permutes \((a,b,c,d)\) to \((c,d,a,b)\),
negating \(h\) to \((b,a,d,c)\), and exchanging them to
\((a,c,b,d)\). Substitution in the four criteria proves the
stated extreme/complete transformations; no strong interchange claim
for the two irrelevance forms is needed.

Keynes's strict condition forbids every newly entailed relevant
content-part. The complete-irrelevance criterion therefore gives
\(e\models i\) unless the prior decides \(h\); compatibility
excludes Carnap's impossible-evidence case. In the balanced binary
example, the singular predictive formula of the preceding Carnap
chapter gives the three fractions displayed in the reader and their
strict signs. Both additional observations together add one positive
and one negative, restoring one half. This verifies cancellation rather
than irrelevance of every component. Tautologies and context-entailed
sentences have relevance zero, explaining why a definition requiring
every weaker part to be positive would be unsuitable. The purely
logical \(C'\) criterion in §86 follows as well: it says
\(a,d>0\) and \(b=0\) or \(c=0\).
\end{proof}
\end{annexitem}

\begin{annexitem}{annex:p15.carnap-infinite-relevance-boundary}{thm:p15-carnap-infinite-relevance-boundary}{The original null-limit countercases}
\begin{proof}
The binary universal-law formula already proved in the preceding
Carnap chapter gives \((s+1)/(N+1)\) and, after one further
positive, \((s+2)/(N+1)\). Their limits are zero although each
finite difference is positive. The singular prediction is
\((s+1)/(s+2)\), whereas the universal hypothesis entails it,
giving conditional value one in every finite system. This proves the
claimed asymmetry of the limiting relevance notions.

The sentence \(j\) permits exactly two complete states: all positive,
with mass \(1/(N+1)\), and only \(c\) negative, with mass
\(1/(N(N+1))\). Therefore \(m_N^*(j)=1/N\).
It implies \(e\), while \(h\) implies \(j\), so
\(c_N^*(h,ej)=m_N^*(h)/m_N^*(j)=N/(N+1)\).
Since \(m_N^*(e)=1/(s+1)\), the signed measure is
\[
 r_N^*(j,h,e)=\frac1{N+1}\left(\frac1{s+1}-\frac1N\right),
\]
which tends to zero. The difference of the conditional limits is one.
The limiting denominator \(m^*(ej)\) is zero; replacing a limit
of quotients with the quotient of their limits is invalid here.
These are explicit reconstructions of the exceptions discussed after
T65-6--8 and T67-9 \cite{Carnap1950Probability}.
\end{proof}
\end{annexitem}
\begin{annexitem}{annex:p15.carnap-general-mean}{thm:p15-carnap-general-mean}{Conditional distributions and the general integral estimate}
\begin{proof}
For \(P(E)>0\), intersection with \(E\) preserves disjoint unions;
countable additivity and normalization of \(P\) therefore give a
probability measure \(P_E\). For real measurable \(X\), define
\(\mu_e(B)=P_E(X^{-1}(B))\) on Borel subsets of the real line.
Preimages preserve disjoint unions and \(X^{-1}(\mathbb R)=\Omega\),
so this is again a probability measure.
For a nonnegative simple Borel function \(g\), summing its values
over its measurable level sets gives
\(\int g(X)\,dP_E=\int g(x)\,d\mu_e(x)\).
Increasing simple approximations and the already proved monotone
convergence theorem, Theorem~\ref{thm:p15-monotone-integral}, extend
this to nonnegative measurable \(g\). Apply it to positive and
negative parts for integrable signed functions. In particular the
finite support law gives exactly \(\sum_j v_jc(h_j,e)\).
The hypotheses are exhaustive, so their probabilities sum to one;
there is no additional normalizing denominator.

Theorem~\ref{thm:p15-signed-integral-laws} proves integrability and
linearity of every real finite linear combination. Apply it to
\(P_E\), whose integral of the constant one is one, to obtain the
reader's affine estimate identity. The interval identity follows from
the disjoint decomposition \(({-\infty},b]=({-\infty},a]\cup(a,b]\).
The Stieltjes notation names integration against this already specified
law, rather than asserting the existence of a law from unverified
sentence assignments. If it has density \(p_e\), the integral
identity follows first for indicators, then for simple functions and
by monotone convergence and signed parts as above.

For the stated counterexample, integrating
\(1/(\pi(1+x^2))\) over \(\mathbb R\) gives one from the
antiderivative \(\arctan(x)/\pi\). Thus it defines a probability
density. But
\(\int_0^R x/(\pi(1+x^2))\,dx=
\log(1+R^2)/(2\pi)\to\infty\); symmetry gives the same
divergence for the negative part. Its signed mean is consequently
undefined, not zero by cancellation. The lottery's weighted sum is
\(10(0.01)+2(0.15)=0.4\); its possible values and its mode are
as stated. The two equal-weight three-value means are respectively
\((3+4+5)/3=4\) and \((3+4+500)/3=169\), while their middle
ordered value remains four. The source's definition and general-language
qualification are D100-1 and pp.~520--528 \cite{Carnap1950Probability}.
\end{proof}
\end{annexitem}

\begin{annexitem}{annex:p15.carnap-estimation-reliability}{thm:p15-carnap-estimation-reliability}{Signed error, square loss, and reliability countercases}
\begin{proof}
Write \(\mu=\int X\,dP_E\). Linearity gives
\(\int(\mu-X)\,dP_E=0\). Square integrability implies
integrability because \(|X|\leq1+X^2\). Expanding the square and
using linearity gives
\[
 \int(X-a)^2\,dP_E=\int X^2\,dP_E-2a\mu+a^2
 =\bigl(\int X^2\,dP_E-\mu^2\bigr)+(a-\mu)^2.
\]
For \(a=\mu\) this proves the error identity and nonnegativity;
the last square proves the unique minimum at \(a=\mu\).
The zero-integral criterion for nonnegative functions,
Theorem~\ref{thm:p15-nonnegative-integral-laws}, says the error is
zero exactly when \((X-\mu)^2=0\) almost surely. In a finite
regular model every state satisfying the evidence has positive
conditional mass, so almost-sure constancy is constancy on that range.

For values \(1,2,3\), the equiprobable mean is two and mean square
\((1+4+9)/3=14/3\). Deviations from two have squares \(1,0,1\),
so the three error means are respectively \(2/3,2/5,4/5\).
The lottery has second moment \(100(0.01)+4(0.15)=1.6\);
subtracting \(0.4^2\) gives \(1.44\), whose square root is
\(1.2\). For the tail example, symmetry gives mean zero, its
one-unit coverage is \(0.99\), and its second moment is
\(2(0.005)100^2=100\). In the two-point interval example the mean
is one. A centered closed interval of radius less than one contains
neither value; a radius from one through values less than ninety-nine
contains only zero and has probability \(0.99\); radius at least
ninety-nine has probability one. None has probability one half.
These computations establish the displayed countercases as well as
T102-1 and T103-1--2's general mathematical content
\cite{Carnap1950Probability}. Optimality for squared loss is not
optimality for every loss or an external validation of the probability law.
\end{proof}
\end{annexitem}

\begin{annexitem}{annex:p15.carnap-frequency-reliability}{thm:p15-carnap-frequency-reliability}{Dependent truth counts and both frequency-error regimes}
\begin{proof}
Put \(I_j=\mathbf1_{R(h_j)}\). Their sum is the truth count,
including logically equivalent but distinct expressions. Linearity of
the finite conditional expectation gives \(\widehat T_e=\sum_jc(h_j,e)\);
division by the positive list size gives its relative-frequency mean.
No independence enters.
For the general truth-fraction error, expand
\(\operatorname{Var}_E(\sum I_j)\) into the sum of its individual
variances and twice its distinct-pair covariances. Indicator variance
is \(p_j(1-p_j)\); pair covariance is
\(c(h_jh_l,e)-p_jp_l\). Dividing by \(t^2\) gives the
general formula, including equivalent or overlapping hypotheses.
Under individual symmetry, permutations of new
For the general truth-fraction error, expand
\(\operatorname{Var}_E(\sum I_j)\) into the sum of its individual
variances and twice its distinct-pair covariances. Indicator variance
is \(p_j(1-p_j)\); pair covariance is
\(c(h_jh_l,e)-p_jp_l\). Dividing by \(t^2\) gives the
general formula, including equivalent or overlapping hypotheses.
Under individual symmetry, permutations of new
names leave the evidence fixed and equate the singular prediction
values. Their arithmetic mean is therefore the common value \(p\).

For direct sampling with known counts, individual symmetry makes every
specified population assignment equally probable on that evidence, as
proved in the preceding Carnap chapter's direct-inference theorem.
Consequently \(E_E I_j=K/N=p\), and for distinct sample positions
\(E_E(I_jI_l)=K(K-1)/(N(N-1))\). Subtracting \(p^2\) gives
covariance \(-p(1-p)/(N-1)\). Expanding the square of the sum
of \(t\) indicators proves
\[
 \operatorname{Var}_E(S)=tp(1-p)-t(t-1)p(1-p)/(N-1)
 =tp(1-p)(N-t)/(N-1).
\]
Scaling by \(t^{-2}\) gives the relative-frequency error; adding
the square of the mean gives the second moment. When \(N=1\)
the only sample is determined. As \(N\to\infty\) with
\(t/N\to0\), \((N-t)/(N-1)\to1\), proving the stated
approximation qualification.

The hypergeometric mass at feasible count \(u\) is
\(f(u)=\binom Ku\binom{N-K}{t-u}/\binom Nt\).
For successive feasible counts,
\[
 \frac{f(u)}{f(u-1)}=
 \frac{(K-u+1)(t-u+1)}{u(N-K-t+u)}.
\]
Its positive denominator allows comparison with one; subtracting
denominator from numerator gives
\((K+1)(t+1)-u(N+2)\). The mass increases until \(u=q\)
and decreases thereafter, with equality at an integer \(q\).
The feasible lower and upper limits are
\(L=\max(0,t+K-N)\), \(U=\min(K,t)\).
We have \(q<U+1\), and \(q>L\): for positive \(L\),
\(q-L=(N-t+1)(N-K+1)/(N+2)>0\); for \(L=0\) it is
immediate. Hence the specified modes are feasible and exhaustive,
including degenerate endpoints.

For predictive indicators, each has variance \(p(1-p)\), and
each distinct-pair covariance is \(\gamma=q_2-p^2\).
The same expansion yields
\(\operatorname{Var}_E(S/t)=p(1-p)/t+(t-1)\gamma/t\).
If these pair marginals fit arbitrarily large exchangeable samples,
nonnegativity of every such variance forces \(\gamma\geq0\)
by taking their lengths to infinity. The displayed variance then
tends to \(\gamma\); this proves a numerical limit without
asserting an actual frequency limit. For binary \(c^*\), the
already proved predictive block formula gives
\(q_2=(s_++1)(s_++2)/((s+2)(s+3))\).
Subtracting \(p^2\) simplifies to \(p(1-p)/(s+3)\).
For \(s=0\) and for \(s=100,s_+=50\) it gives
\(1/12,1/412\). At \(N=14,K=10,t=7\), substitution in
the exact direct formula gives \(10/13\), whereas its uncorrected
version is \(10/7\). Adding the observed count to the predictive
mean of the remainder gives the inverse mean in the reader.

All these proofs concern the general results of §§104--107 or the
explicitly selected \(c^*\) cases, not a later continuum. In
particular the constancy of a sequence of estimated relative frequencies
does not prove convergence of the actual values of a given infinite
class. Repeated positive and negative blocks of lengths increasing so
fast that each dominates all previous blocks give an enumeration with
no frequency limit; alternating positives and negatives give another
enumeration of the same countably infinite two-kind class with limit
one half. This verifies the order/existence distinction in §106B.
\end{proof}
\end{annexitem}
